\chapter{Power Series and Taylor Series}

Somewhere in calculus the pieces you have collected fall into a single picture. Polynomials are the friendliest functions of all: you can add them, multiply them, differentiate them, integrate them, evaluate them with nothing but arithmetic. The transcendental functions $e^x$, $\sin x$, $\cos x$, $\ln(1+x)$ seemed to belong to a different species. You could not write $\sin x$ as any finite combination of powers of $x$, so to compute $\sin(0.3)$ you reached for a calculator and trusted a black box.

This chapter dissolves that wall. \emph{Nearly every function you care about is secretly a polynomial -- one of infinite degree.} The function $e^x$, the function $\sin x$, the function $\ln(1+x)$ -- each is equal, on some interval, to a genuine sum of powers of $x$,
\[
  c_0 + c_1 x + c_2 x^2 + c_3 x^3 + \cdots,
\]
where the coefficients $c_n$ are specific, computable numbers. Once you know those coefficients, the mysterious transcendental function becomes something you can compute by hand, differentiate term by term, integrate term by term, and reason about with the comfortable algebra of polynomials.

These infinite polynomials are called \emph{power series}, and the recipe that produces the coefficients is the \emph{Taylor series}. By the end of this chapter you will see not only \emph{how} to build them but \emph{why} the recipe must give the coefficients it does -- it is forced, with no freedom at all, by demanding that the polynomial match the function's value, slope, curvature, and every higher rate of change at a single point.

\section{What a Power Series Is}
\label{ch30:sec-what}

\begin{intuitionbox}
A power series is a polynomial that never stops. Think of it as a recipe for a function written entirely in powers of $x$: pick a list of numbers $c_0, c_1, c_2, \dots$, and the series is
\[
  c_0 + c_1 x + c_2 x^2 + c_3 x^3 + \cdots
\]
For most lists this sum converges only for $x$ close to $0$, and there it defines a genuine function. The family of functions you can reach this way turns out to be remarkably rich.
\end{intuitionbox}

\begin{definition}[Power series]
\label{ch30:def-power-series}
A \emph{power series centered at} $a$ is an expression of the form
\[
  \sum_{n=0}^{\infty} c_n (x-a)^n
  \;=\; c_0 + c_1(x-a) + c_2(x-a)^2 + c_3(x-a)^3 + \cdots,
\]
where the numbers $c_n$ are the \emph{coefficients} and the fixed number $a$ is the \emph{center}. For each fixed value of $x$ this is an infinite series of real numbers in the sense of the previous chapter, and we ask the same question we always ask of an infinite sum: does it converge?
\end{definition}

When $a=0$ the series simplifies to $\sum c_n x^n$, and that is the case we will live in most of the time. A series centered at $a$ is nothing more than the same idea shifted so that its natural ``home base'' sits at $x=a$ instead of at the origin. Everything we prove about $\sum c_n x^n$ near $0$ transfers, by the substitution $x \mapsto x-a$, to $\sum c_n (x-a)^n$ near $a$.

A power series is a \emph{function of $x$}, but only on the set of $x$ where the series actually converges. At $x=a$ every term past the first vanishes, so the series collapses to $c_0$ -- it always converges at its center. The interesting question is how far away from the center the convergence survives.

\begin{example}[The geometric series is a power series]
\label{ch30:ex-geometric}
The most important power series in all of mathematics is one you have already met:
\[
  1 + x + x^2 + x^3 + \cdots = \sum_{n=0}^{\infty} x^n.
\]
Here every coefficient is $c_n = 1$. From the previous chapter you know this geometric series converges precisely when $\abs{x} < 1$, and when it does, it sums to $\dfrac{1}{1-x}$. So we may write the exact equation
\[
  \frac{1}{1-x} = \sum_{n=0}^{\infty} x^n, \qquad \abs{x} < 1.
\]
This is our first function-as-infinite-polynomial. The rational function $1/(1-x)$, defined everywhere except $x=1$, agrees with the infinite polynomial on the open interval $(-1,1)$ and disagrees (the series diverges) everywhere else. Hold onto this example: almost every other series in this chapter is built from it by differentiating, integrating, or substituting.
\end{example}

\subsection{Where does a power series live? Radius of convergence}

For the geometric series the set of good $x$ was the symmetric interval $(-1,1)$. That symmetry is no accident. The deep structural fact about power series is that the set of points where they converge is \emph{always} an interval centered at $a$ -- there is a single radius $R$ such that the series converges inside and diverges outside.

\begin{theorem}[Radius of convergence]
\label{ch30:thm-radius}
For any power series $\sum_{n=0}^{\infty} c_n (x-a)^n$, exactly one of the following holds:
\begin{enumerate}
  \item The series converges only at $x = a$ (we say $R = 0$).
  \item The series converges for all real $x$ (we say $R = \infty$).
  \item There is a positive number $R$, the \emph{radius of convergence}, such that the series converges absolutely whenever $\abs{x-a} < R$ and diverges whenever $\abs{x-a} > R$.
\end{enumerate}
\end{theorem}

\begin{intuitionbox}
Picture a circle of influence around the center $a$. Inside the circle the powers $(x-a)^n$ shrink fast enough that the series adds up to something finite. Outside, the powers grow and overwhelm any coefficients, so the sum explodes. The boundary -- the two endpoints $a \pm R$ -- is a delicate place that must be checked by hand each time, because the series can go either way there.
\end{intuitionbox}

The engine behind this theorem is a single comparison idea, resting on one lemma: if a power series converges \emph{anywhere} away from the center, it converges \emph{everywhere closer in}.

\begin{lemma}[Abel's comparison]
\label{ch30:lem-abel}
Suppose $\sum c_n b^n$ converges for some particular $b \neq 0$. Then $\sum c_n x^n$ converges absolutely for every $x$ with $\abs{x} < \abs{b}$.
\end{lemma}

\begin{proof}
If $\sum c_n b^n$ converges, then its terms must march to zero (the term test, from the previous chapter). A sequence that tends to zero is in particular bounded, so there is a constant $M$ with
\[
  \abs{c_n b^n} \le M \qquad \text{for every } n.
\]
Now take any $x$ closer to the origin than $b$, and let $r = \abs{x}/\abs{b}$, a number strictly less than $1$. Watch how the terms of $\sum c_n x^n$ get controlled:
\[
  \abs{c_n x^n}
  = \abs{c_n b^n}\cdot \left|\frac{x}{b}\right|^n
  \le M \, r^n.
\]
We have bounded our terms by the terms of a convergent geometric series $\sum M r^n$ (convergent because $r<1$). By the comparison test, $\sum \abs{c_n x^n}$ converges. That is exactly absolute convergence for $\sum c_n x^n$.
\end{proof}

\begin{whybox}
Why does the convergence set have to be an \emph{interval} -- why couldn't a power series converge at $x=3$ and $x=7$ but diverge at $x=5$ in between? Abel's lemma forbids it. Convergence at the far point $x=7$ would force convergence at every point closer to the center than $7$, including $x=5$. So the ``good'' set has no holes: if it contains a point, it contains everything between that point and the center. A subset of the line that contains the whole segment from the center out to each of its points, and is symmetric about the center, can only be a single interval centered at $a$. That is the entire content of Theorem~\ref{ch30:thm-radius}: the radius $R$ is simply the supremum of the distances $\abs{b-a}$ at which convergence happens.
\end{whybox}

\begin{figure}[H]
\centering
\begin{tikzpicture}
  \draw[->] (-6,0) -- (6,0) node[right] {$x$};
  % center
  \filldraw[structurepurple] (0,0) circle (2pt) node[below=5pt] {$a$};
  % converge region
  \draw[curveblue, line width=3pt] (-4,0) -- (4,0);
  \node[curveblue] at (0,0.55) {converges (absolutely)};
  % endpoints
  \filldraw[white,draw=trapred,line width=1.2pt] (-4,0) circle (2.6pt);
  \filldraw[white,draw=trapred,line width=1.2pt] (4,0) circle (2.6pt);
  \node[trapred] at (-4,-0.5) {$a-R$};
  \node[trapred] at (4,-0.5) {$a+R$};
  \node[trapred] at (4,0.45) {\small ?};
  \node[trapred] at (-4,0.45) {\small ?};
  % diverge
  \draw[trapred, line width=1.5pt, dashed] (4.05,0) -- (5.8,0);
  \draw[trapred, line width=1.5pt, dashed] (-5.8,0) -- (-4.05,0);
  \node[trapred] at (5.1,0.45) {diverges};
  \node[trapred] at (-5.1,0.45) {diverges};
  % radius brace
  \draw[decorate,decoration={brace,amplitude=6pt},gray] (0,-0.9) -- (4,-0.9);
  \node[gray] at (2,-1.4) {radius $R$};
\end{tikzpicture}
\caption{The interval of convergence. Inside the radius $R$ the series converges absolutely; outside it diverges. The two endpoints are question marks that must be tested individually.}
\label{ch30:fig-interval}
\end{figure}

\subsection{Finding the radius with the ratio test}

In practice you almost always find $R$ by feeding the series to the ratio test from the previous chapter. The ratio test asks for the limit of the absolute ratio of consecutive terms; convergence happens when that limit is below $1$. For a power series the variable $x$ rides along in that ratio, and the condition ``ratio $< 1$'' becomes a condition on $\abs{x}$ -- which is exactly the radius.

\begin{focusbox}
\textbf{The ratio-test recipe for the radius.} Apply the ratio test to $\sum c_n(x-a)^n$:
\[
  L(x) = \lim_{n\to\infty}\left|\frac{c_{n+1}(x-a)^{n+1}}{c_n(x-a)^n}\right|
       = \abs{x-a}\cdot \lim_{n\to\infty}\left|\frac{c_{n+1}}{c_n}\right|.
\]
If $\lim \abs{c_{n+1}/c_n} = \ell$, then convergence ($L<1$) means $\abs{x-a} < 1/\ell$, so
\[
  \boxed{\,R = \frac{1}{\ell} = \lim_{n\to\infty}\left|\frac{c_n}{c_{n+1}}\right|\,}
\]
(with $R=\infty$ if $\ell=0$ and $R=0$ if $\ell=\infty$). Always remember to test the endpoints $a\pm R$ separately.
\end{focusbox}

\begin{example}[Radius of the exponential series]
\label{ch30:ex-exp-radius}
Consider $\displaystyle\sum_{n=0}^{\infty}\frac{x^n}{n!}$, so $c_n = 1/n!$. The ratio of consecutive coefficients is
\[
  \left|\frac{c_{n+1}}{c_n}\right| = \frac{1/(n+1)!}{1/n!} = \frac{n!}{(n+1)!} = \frac{1}{n+1} \To 0.
\]
So $\ell = 0$ and $R = 1/\ell = \infty$. This series converges for \emph{every} real $x$. The factorials in the denominator grow so violently fast that they crush any value of $x^n$ in the numerator, no matter how large $x$ is. (We will soon see this series is $e^x$.)
\end{example}

\begin{example}[A series with a finite radius and interesting endpoints]
\label{ch30:ex-ln-radius}
Consider $\displaystyle\sum_{n=1}^{\infty}\frac{x^n}{n}$, so $c_n = 1/n$. Then
\[
  \left|\frac{c_n}{c_{n+1}}\right| = \frac{1/n}{1/(n+1)} = \frac{n+1}{n} \To 1,
\]
so $R = 1$. The series converges on $(-1,1)$ and diverges outside $[-1,1]$. Now the endpoints:
\begin{itemize}
  \item At $x = 1$ the series is $\sum 1/n$, the harmonic series -- diverges.
  \item At $x = -1$ the series is $\sum (-1)^n/n$, the alternating harmonic series -- converges.
\end{itemize}
So the true interval of convergence is $[-1, 1)$: closed on the left, open on the right. This lopsidedness is exactly why the endpoints cannot be guessed and must be checked. (We will soon see this series is $-\ln(1-x)$.)
\end{example}

\begin{labbox}
Find the radius of convergence of $\displaystyle\sum_{n=0}^{\infty} n!\, x^n$ by hand. Compute $\abs{c_n/c_{n+1}} = n!/(n+1)! = 1/(n+1) \to 0$, so $R = 0$: this series converges \emph{only} at $x=0$. The coefficients grow so fast that no nonzero $x$ can tame them. A power series really can have a vanishing radius.
\end{labbox}

\begin{trapbox}
The radius is determined by the coefficients, \emph{not} by where the function you are representing happens to misbehave -- but the two are often related. Be careful, though, with two classic confusions:
\begin{itemize}
  \item ``$R$ large means the terms are large.'' No: large $R$ means the coefficients $c_n$ \emph{shrink} fast, so the powers can grow far before overwhelming them.
  \item ``Convergence at one endpoint guarantees the other.'' No, as Example~\ref{ch30:ex-ln-radius} shows. Each endpoint is its own separate question.
\end{itemize}
\end{trapbox}

\section{Building Functions from the Geometric Series}
\label{ch30:sec-building}

Before we develop the general Taylor machinery, let us see how much mileage we get from the single equation
\[
  \frac{1}{1-x} = \sum_{n=0}^{\infty} x^n, \qquad \abs{x} < 1,
\]
together with two tools we are about to justify: \emph{you may differentiate and integrate a power series term by term inside its radius of convergence}, and \emph{you may substitute one power series into a simple expression}. These manipulations are the bread and butter of the subject, and they let us grow a whole orchard of series from one seed.

\begin{theorem}[Term-by-term differentiation and integration]
\label{ch30:thm-termwise}
Suppose $f(x) = \sum_{n=0}^{\infty} c_n (x-a)^n$ has radius of convergence $R > 0$. Then on the open interval $(a-R, a+R)$ the function $f$ is differentiable (indeed infinitely differentiable), and its derivative and antiderivative are obtained term by term:
\[
  f'(x) = \sum_{n=1}^{\infty} n\, c_n (x-a)^{n-1},
  \qquad
  \int f(x)\,\dd x = C + \sum_{n=0}^{\infty} \frac{c_n}{n+1}(x-a)^{n+1}.
\]
Both new series have the same radius of convergence $R$.
\end{theorem}

We will not prove this in full -- the careful proof belongs to the theory of uniform convergence -- but the next box explains why it ought to be true, and why the radius is preserved.

\begin{whybox}
A finite polynomial can of course be differentiated term by term; that is just the power rule applied to each piece. The only thing in doubt for an \emph{infinite} sum is whether the infinitely many error contributions, when you swap ``differentiate'' and ``add up'', stay under control. They do, and the reason is the absolute convergence guaranteed \emph{strictly inside} the radius (Lemma~\ref{ch30:lem-abel}). Inside the radius there is room to spare -- the terms decay geometrically -- and that slack is exactly what lets the limit defining the derivative pass through the infinite sum. As for the radius being unchanged: differentiating multiplies the coefficient $c_n$ by $n$, and $\lim \abs{n c_n / ((n+1)c_{n+1})} = \lim (n/(n+1))\cdot\abs{c_n/c_{n+1}} = R$, since $n/(n+1)\to 1$. The factor of $n$ is too mild to change where the series lives.
\end{whybox}

\begin{example}[A series for $\dfrac{1}{(1-x)^2}$ by differentiating]
\label{ch30:ex-deriv-geom}
Differentiate both sides of $\dfrac{1}{1-x} = \sum_{n=0}^{\infty} x^n$. The left side differentiates to $\dfrac{1}{(1-x)^2}$. The right side differentiates term by term:
\[
  \ddx{\sum_{n=0}^{\infty} x^n} = \sum_{n=1}^{\infty} n x^{n-1} = 1 + 2x + 3x^2 + 4x^3 + \cdots.
\]
So, for $\abs{x}<1$,
\[
  \frac{1}{(1-x)^2} = \sum_{n=1}^{\infty} n\,x^{n-1} = \sum_{n=0}^{\infty}(n+1)x^n.
\]
We obtained a brand-new, nontrivial series for free, just by differentiating the one series we knew.
\end{example}

\begin{example}[The series for $\ln(1+x)$ by integrating]
\label{ch30:ex-ln-series}
Start from the geometric series with $x$ replaced by $-t$:
\[
  \frac{1}{1+t} = \frac{1}{1-(-t)} = \sum_{n=0}^{\infty}(-t)^n = 1 - t + t^2 - t^3 + \cdots, \qquad \abs{t}<1.
\]
Now integrate both sides from $0$ to $x$. On the left,
\[
  \int_0^x \frac{\dd t}{1+t} = \ln(1+x).
\]
On the right, integrate term by term:
\[
  \int_0^x \left(1 - t + t^2 - t^3 + \cdots\right)\dd t
  = x - \frac{x^2}{2} + \frac{x^3}{3} - \frac{x^4}{4} + \cdots.
\]
Therefore
\[
  \boxed{\,\ln(1+x) = \sum_{n=1}^{\infty} \frac{(-1)^{n-1}}{n}\,x^n
  = x - \frac{x^2}{2} + \frac{x^3}{3} - \frac{x^4}{4} + \cdots\,}, \qquad -1 < x \le 1.
\]
Notice the surprise at the endpoint: although we derived the formula assuming $\abs{x}<1$, at $x=1$ the right side becomes the convergent alternating harmonic series, and one can show it does equal $\ln 2$. So
\[
  \ln 2 = 1 - \frac12 + \frac13 - \frac14 + \frac15 - \cdots,
\]
a strange and beautiful identity: an alternating sum of unit fractions equals the natural logarithm of $2$.
\end{example}

\begin{example}[The series for $\arctan x$, and a formula for $\pi$]
\label{ch30:ex-arctan}
The same trick produces one of the most famous series in history. Replace $x$ by $-t^2$ in the geometric series:
\[
  \frac{1}{1+t^2} = \sum_{n=0}^{\infty}(-t^2)^n = 1 - t^2 + t^4 - t^6 + \cdots, \qquad \abs{t}<1.
\]
Integrate from $0$ to $x$, using $\int_0^x \dfrac{\dd t}{1+t^2} = \arctan x$:
\[
  \boxed{\,\arctan x = x - \frac{x^3}{3} + \frac{x^5}{5} - \frac{x^7}{7} + \cdots
  = \sum_{n=0}^{\infty}\frac{(-1)^n}{2n+1}\,x^{2n+1}\,}, \qquad -1 \le x \le 1.
\]
Setting $x=1$ gives Leibniz's stunning formula:
\[
  \frac{\pi}{4} = 1 - \frac13 + \frac15 - \frac17 + \frac19 - \cdots.
\]
\end{example}

\begin{historybox}
The series for $\arctan x$ -- and with it the formula above for $\pi/4$ -- was discovered in the fourteenth century by the Indian mathematician \textbf{Madhava of Sangamagrama} and his Kerala school, roughly three hundred years before Newton and Leibniz. They had series for sine, cosine, and arctangent, and used them to compute $\pi$ to eleven decimal places. In Europe the arctangent series is usually credited to James Gregory (1671) and Leibniz (1673). It is a humbling reminder that the ``infinite polynomial'' idea was glimpsed, with real rigor for its time, far from the usual narrative centers. The Leibniz series converges agonizingly slowly -- you need hundreds of terms for a couple of digits -- but it was the first time a simple, exact, infinite arithmetic recipe was given for the circle's most famous constant.
\end{historybox}

\section{Taylor and Maclaurin Series: the Universal Recipe}
\label{ch30:sec-taylor}

So far we have manufactured series by clever manipulation of the geometric series. That is wonderful when it works, but it does not work for $e^x$, $\sin x$, or $\cos x$ -- those are not built from $1/(1-x)$ by any obvious substitution. We need a \emph{universal} method that takes \emph{any} sufficiently smooth function and produces its power series directly. That method is Taylor's, and the idea behind it is one of the cleanest in all of calculus.

\begin{intuitionbox}
Suppose a function $f$ really is equal to some power series near a point $a$:
\[
  f(x) = c_0 + c_1(x-a) + c_2(x-a)^2 + c_3(x-a)^3 + \cdots.
\]
We do not yet know the coefficients. But we have a tool that extracts them one at a time: \emph{differentiation}. Each derivative, evaluated at the center $a$, reaches in and pulls out exactly one coefficient, leaving the others either already gone or not yet arrived. The coefficients are not free; they are \emph{forced} to be the derivatives of $f$ at $a$.
\end{intuitionbox}

Let us watch the extraction happen. Assume $f(x) = \sum_{n=0}^\infty c_n (x-a)^n$ and that we may differentiate term by term.

First, just plug in $x = a$. Every term with a factor of $(x-a)$ dies, leaving only the constant:
\[
  f(a) = c_0.
\]
Next, differentiate once: $f'(x) = c_1 + 2c_2(x-a) + 3c_3(x-a)^2 + \cdots$, then plug in $x=a$ again. Once more all the $(x-a)$ factors vanish:
\[
  f'(a) = c_1.
\]
Differentiate again: $f''(x) = 2c_2 + 6c_3(x-a) + 12 c_4 (x-a)^2 + \cdots$, and set $x=a$:
\[
  f''(a) = 2c_2 \quad\Longrightarrow\quad c_2 = \frac{f''(a)}{2}.
\]
Once more: $f'''(x) = 6 c_3 + 24 c_4 (x-a) + \cdots$, so $f'''(a) = 6 c_3$, giving $c_3 = f'''(a)/6$. The pattern is unmistakable. Differentiating $n$ times and evaluating at $a$ leaves only the $n$th term, which has been multiplied by $n\cdot(n-1)\cdots 2\cdot 1 = n!$. Hence
\[
  f^{(n)}(a) = n!\, c_n \qquad\Longrightarrow\qquad c_n = \frac{f^{(n)}(a)}{n!}.
\]

\begin{definition}[Taylor and Maclaurin series]
\label{ch30:def-taylor}
Let $f$ be infinitely differentiable at $a$. The \emph{Taylor series of $f$ centered at $a$} is
\[
  \sum_{n=0}^{\infty} \frac{f^{(n)}(a)}{n!}\,(x-a)^n
  = f(a) + f'(a)(x-a) + \frac{f''(a)}{2!}(x-a)^2 + \frac{f'''(a)}{3!}(x-a)^3 + \cdots.
\]
When the center is $a=0$ this is also called the \emph{Maclaurin series}:
\[
  f(0) + f'(0)\,x + \frac{f''(0)}{2!}x^2 + \frac{f'''(0)}{3!}x^3 + \cdots.
\]
The polynomial obtained by stopping at the term of degree $N$ is the \emph{Taylor polynomial} $P_N(x)$ of order $N$.
\end{definition}

\begin{whybox}
Why this exact recipe and no other? Because a power series matching $f$ has \emph{no freedom}. Each coefficient is pinned down by one derivative: $c_n$ controls (and is the only coefficient controlling) the $n$th derivative at the center. The factor $n!$ appears because differentiating $(x-a)^n$ a total of $n$ times produces the falling product $n(n-1)\cdots 1 = n!$. The Taylor recipe is simply the unique answer to the question: \emph{which polynomial has the same value, the same slope, the same curvature, the same third derivative, and so on, all the way up, as $f$ at the point $a$?} You are not approximating by guesswork; you are matching $f$ to infinite order at a single point.
\end{whybox}

\begin{happeningbox}
Here is what a Taylor polynomial \emph{really} is, stripped of symbols. The order-$0$ polynomial $P_0 = f(a)$ is the horizontal line through the point -- it gets the height right. The order-$1$ polynomial $P_1 = f(a) + f'(a)(x-a)$ is the tangent line -- it gets the height \emph{and} the slope right; this is exactly the linear approximation from the chapter on differentials. The order-$2$ polynomial adds $\tfrac{1}{2}f''(a)(x-a)^2$, a parabola that also matches the \emph{curvature}. Each new term fixes the next derivative without disturbing the ones already matched, because $(x-a)^{n}$ and all its first $n-1$ derivatives vanish at $a$. So the Taylor series is a tower: each floor corrects a finer feature of the function's shape at the center, and ``infinitely many floors'' means matching the function's local shape to infinite precision. Whether that local matching extends to a global equality is the subtle question we take up in the next chapter.
\end{happeningbox}

\subsection{The three great Maclaurin series}

We now compute, from the definition, the three series that every scientist carries in their head: $e^x$, $\sin x$, $\cos x$. Each is centered at $0$, and in each case the derivatives at $0$ fall into a simple repeating pattern.

\begin{example}[The series for $e^x$]
\label{ch30:ex-exp-series}
Let $f(x) = e^x$. The defining magic of the exponential is that it is its own derivative: $f'(x) = e^x$, $f''(x) = e^x$, and in general $f^{(n)}(x) = e^x$ for all $n$. Evaluated at $0$, every derivative is $e^0 = 1$:
\[
  f^{(n)}(0) = 1 \qquad \text{for all } n.
\]
So the coefficients are $c_n = f^{(n)}(0)/n! = 1/n!$, and
\[
  \boxed{\,e^x = \sum_{n=0}^{\infty} \frac{x^n}{n!}
  = 1 + x + \frac{x^2}{2!} + \frac{x^3}{3!} + \frac{x^4}{4!} + \cdots\,}, \qquad \text{all } x.
\]
We computed the radius of this series back in Example~\ref{ch30:ex-exp-radius}: it is infinite. The same factorial growth that gave $R=\infty$ now reveals itself as the engine of $e^x$.
\end{example}

\begin{happeningbox}
This series is arguably the \emph{true definition} of $e^x$. In the chapter on the exponential we met $e^x$ three ways: as a limit of compound interest, as the function equal to its own derivative, and as the inverse of $\ln$. The series $\sum x^n/n!$ unifies them. It is its own derivative by term-by-term differentiation: differentiating $x^n/n!$ gives $x^{n-1}/(n-1)!$, which is just the previous term, so the whole series reproduces itself. It equals $1$ at $x=0$. Those two facts -- self-derivative and value $1$ at $0$ -- determine $e^x$ completely. Setting $x=1$ even gives the number $e$ itself as a rapidly converging sum: $e = 1 + 1 + \tfrac12 + \tfrac16 + \tfrac1{24} + \cdots = 2.71828\ldots$, far faster than the compound-interest limit ever did.
\end{happeningbox}

\begin{example}[The series for $\sin x$ and $\cos x$]
\label{ch30:ex-sin-cos}
Let $f(x) = \sin x$. Its derivatives cycle with period four:
\[
  f = \sin x,\quad f' = \cos x,\quad f'' = -\sin x,\quad f''' = -\cos x,\quad f'''' = \sin x,\ \dots
\]
Evaluating at $0$, where $\sin 0 = 0$ and $\cos 0 = 1$, we get the pattern $0, 1, 0, -1, 0, 1, 0, -1, \dots$. The even-order derivatives vanish; the odd ones alternate $+1, -1$. So only odd powers survive:
\[
  \boxed{\,\sin x = x - \frac{x^3}{3!} + \frac{x^5}{5!} - \frac{x^7}{7!} + \cdots
  = \sum_{n=0}^{\infty}\frac{(-1)^n}{(2n+1)!}\,x^{2n+1}\,}, \qquad \text{all } x.
\]
The same computation for $\cos x$ (whose derivatives at $0$ run $1, 0, -1, 0, 1, \dots$) keeps only even powers:
\[
  \boxed{\,\cos x = 1 - \frac{x^2}{2!} + \frac{x^4}{4!} - \frac{x^6}{6!} + \cdots
  = \sum_{n=0}^{\infty}\frac{(-1)^n}{(2n)!}\,x^{2n}\,}, \qquad \text{all } x.
\]
Both have infinite radius. Notice that $\sin$ uses only odd powers (it is an odd function) and $\cos$ only even powers (it is even): the series \emph{remembers} the symmetry of the function. And differentiating the $\sin$ series term by term gives the $\cos$ series exactly, just as it must.
\end{example}

\begin{connectionbox}
Line up the three series and a miracle leaps out:
\[
  e^x = 1 + x + \tfrac{x^2}{2!} + \tfrac{x^3}{3!} + \cdots,\quad
  \cos x = 1 - \tfrac{x^2}{2!} + \tfrac{x^4}{4!} - \cdots,\quad
  \sin x = x - \tfrac{x^3}{3!} + \tfrac{x^5}{5!} - \cdots.
\]
The terms of $e^x$ are precisely the terms of $\cos x$ and $\sin x$ interleaved -- if only we could supply the alternating signs. Feed an imaginary argument $ix$ into the exponential series and the powers $i^n$ cycle $1, i, -1, -i, \dots$, producing exactly those signs. The even powers reconstruct $\cos x$, the odd powers reconstruct $i\sin x$, and you land on \textbf{Euler's formula}
\[
  e^{ix} = \cos x + i \sin x,
\]
the deepest identity connecting exponential growth and circular motion. At $x=\pi$ it collapses to $e^{i\pi} + 1 = 0$. The power series is what makes this rigorous: it is the common language in which $e^x$, $\cos x$, and $\sin x$ are visibly the same kind of object.
\end{connectionbox}

\section{Watching the Polynomials Close In}
\label{ch30:sec-watching}

Definitions are one thing; \emph{seeing} a polynomial chase down a transcendental function is another. The Taylor polynomials $P_N(x)$ are ordinary finite polynomials -- you can graph them. As $N$ grows, they hug the true function over a wider and wider stretch. Let us look.

\begin{figure}[H]
\centering
\begin{tikzpicture}
\begin{axis}[
  width=12cm, height=7cm,
  xmin=-7, xmax=7, ymin=-2.2, ymax=2.2,
  axis lines=middle,
  xlabel=$x$, ylabel=$y$,
  samples=200, domain=-7:7,
  legend pos=north west, legend cell align=left,
  restrict y to domain=-2.4:2.4,
]
  \addplot[curveblue, very thick, domain=-7:7] {sin(deg(x))};
  \addlegendentry{$\sin x$}
  \addplot[trapred, thick, domain=-7:7] {x};
  \addlegendentry{$P_1 = x$}
  \addplot[quotientorange, thick, domain=-7:7] {x - x^3/6};
  \addlegendentry{$P_3$}
  \addplot[targetgreen, thick, domain=-7:7] {x - x^3/6 + x^5/120};
  \addlegendentry{$P_5$}
  \addplot[structurepurple, thick, domain=-7:7] {x - x^3/6 + x^5/120 - x^7/5040};
  \addlegendentry{$P_7$}
\end{axis}
\end{tikzpicture}
\caption{Taylor polynomials of $\sin x$ at $0$. $P_1$ (the tangent line) is good only very near the origin; each higher-order polynomial clings to the sine curve over a longer interval before peeling away. By $P_7$ the match is excellent across a full period.}
\label{ch30:fig-sin-taylor}
\end{figure}

The picture tells the whole story. Near the center every polynomial is a good fit -- that is guaranteed, because they match $\sin x$ to high order \emph{at the origin}. The interesting drama is at the edges: a low-order polynomial cannot help shooting off to $\pm\infty$ (a degree-$7$ polynomial must eventually behave like $\pm x^7$), and so it eventually abandons the bounded sine wave. But it holds on longer than its predecessor, and the place where it ``lets go'' marches steadily outward as $N$ increases. Take $N\to\infty$ and the letting-go point runs off to infinity: the full series equals $\sin x$ everywhere.

\begin{figure}[H]
\centering
\begin{tikzpicture}
\begin{axis}[
  width=12cm, height=6.5cm,
  xmin=-3, xmax=3, ymin=-1, ymax=9,
  axis lines=middle,
  xlabel=$x$, ylabel=$y$,
  samples=120, domain=-3:3,
  legend pos=north west, legend cell align=left,
]
  \addplot[curveblue, very thick, domain=-3:3] {exp(x)};
  \addlegendentry{$e^x$}
  \addplot[trapred, thick, domain=-3:3] {1+x};
  \addlegendentry{$P_1$}
  \addplot[quotientorange, thick, domain=-3:3] {1+x+x^2/2};
  \addlegendentry{$P_2$}
  \addplot[targetgreen, thick, domain=-3:3] {1+x+x^2/2+x^3/6};
  \addlegendentry{$P_3$}
\end{axis}
\end{tikzpicture}
\caption{Taylor polynomials of $e^x$ at $0$. The tangent line $P_1=1+x$ falls below the curve immediately; the parabola $P_2$ matches the upward curvature; $P_3$ improves the fit further. Each polynomial sits underneath $e^x$ for $x>0$ and creeps closer as the order rises.}
\label{ch30:fig-exp-taylor}
\end{figure}

\begin{walkbox}
Take a walk and turn this over. The Taylor polynomial matches the function perfectly \emph{at one point} -- value, slope, curvature, and all higher derivatives, to whatever order you choose. Yet a polynomial of degree $7$ must eventually run away to infinity, while $\sin x$ stays politely between $-1$ and $1$ forever. How can these two facts coexist? The resolution: matching to high order at a \emph{point} controls behavior \emph{near} that point, but says nothing directly about far away. The error grows as you leave the center, slowly at first (because so many derivatives agree) and then catastrophically. The whole art of approximation theory is the contest between ``how high the order'' and ``how far from the center'', and the next chapter, on Taylor's theorem with remainder, makes that contest exact.
\end{walkbox}

\begin{labbox}
Compute $\sin(0.5)$ by hand from the series, keeping three terms:
\[
  \sin(0.5) \approx 0.5 - \frac{0.5^3}{6} + \frac{0.5^5}{120}
  = 0.5 - \frac{0.125}{6} + \frac{0.03125}{120}.
\]
That is $0.5 - 0.0208333 + 0.0002604 = 0.4794271$. The true value is $\sin(0.5) = 0.4794255\ldots$, so three terms already give six-digit accuracy. This is literally how a calculator does it: reduce the angle to a small range, then sum a handful of series terms. The black box was a power series all along.
\end{labbox}

\section{Manipulating Series: an Algebra of Infinite Polynomials}
\label{ch30:sec-manipulating}

Inside their radius of convergence, power series obey the same algebra as ordinary polynomials, and that is what makes them so useful. You may add them, subtract them, multiply them, substitute one into another, and (with care) divide them. This means you rarely have to grind out high derivatives from the Taylor definition -- you start from the known series for $e^x$, $\sin x$, $\cos x$, $\ln(1+x)$, $1/(1-x)$ and \emph{build}.

\begin{focusbox}
\textbf{The toolbox of legal moves} (valid wherever the series involved converge):
\begin{itemize}
  \item \textbf{Add / subtract:} combine coefficients of like powers.
  \item \textbf{Multiply by a constant or a power of $x$:} scale or shift every term.
  \item \textbf{Substitute} a simple expression (like $x^2$ or $-x$ or $3x$) for the variable.
  \item \textbf{Multiply two series:} collect terms of each total degree (the Cauchy product).
  \item \textbf{Differentiate / integrate} term by term (Theorem~\ref{ch30:thm-termwise}).
\end{itemize}
The uniqueness of power series (each function has \emph{one} series at a given center) is what makes this legal: whatever route you take, you must arrive at the same coefficients, so you may take the easiest route.
\end{focusbox}

\begin{example}[Substitution: the series for $e^{-x^2}$]
\label{ch30:ex-sub}
The function $e^{-x^2}$ -- the unnormalized Gaussian bell curve -- has no elementary antiderivative, so the integral $\int e^{-x^2}\dd x$ defies the usual techniques. But its series is effortless. Start from $e^u = \sum u^n/n!$ and substitute $u = -x^2$:
\[
  e^{-x^2} = \sum_{n=0}^{\infty}\frac{(-x^2)^n}{n!}
  = \sum_{n=0}^{\infty}\frac{(-1)^n x^{2n}}{n!}
  = 1 - x^2 + \frac{x^4}{2!} - \frac{x^6}{3!} + \cdots.
\]
Now integrate term by term -- something we could never do to the closed form:
\[
  \int_0^x e^{-t^2}\,\dd t
  = x - \frac{x^3}{3} + \frac{x^5}{5\cdot 2!} - \frac{x^7}{7\cdot 3!} + \cdots
  = \sum_{n=0}^{\infty}\frac{(-1)^n x^{2n+1}}{(2n+1)\,n!}.
\]
This is, up to a constant, the error function $\operatorname{erf}$ that governs the normal distribution in statistics. A function with no elementary antiderivative becomes completely computable through its series.
\end{example}

\begin{figure}[H]
\centering
\begin{tikzpicture}
\begin{axis}[
  width=12cm, height=6cm,
  xmin=-3, xmax=3, ymin=0, ymax=1.15,
  axis lines=middle,
  xlabel=$t$, ylabel=$y$,
  samples=120, domain=-3:3,
  ytick={0,0.5,1},
]
  \addplot[areafill!70, draw=none, domain=0:1, samples=60] {exp(-x^2)} \closedcycle;
  \addplot[curveblue, very thick, domain=-3:3] {exp(-x^2)};
  \addlegendentry{$e^{-t^2}$}
  \node[curveblue] at (axis cs:0.45,0.33) {area $=\displaystyle\int_0^{1}\! e^{-t^2}\dd t$};
\end{axis}
\end{tikzpicture}
\caption{The Gaussian bell curve $e^{-t^2}$. The shaded area from $0$ to $1$ has no elementary formula, yet the term-by-term integration of Example~\ref{ch30:ex-sub} turns it into the convergent sum $1 - \tfrac13 + \tfrac{1}{5\cdot 2!} - \tfrac{1}{7\cdot 3!} + \cdots$, which any calculator can add up.}
\label{ch30:fig-gaussian}
\end{figure}

\begin{applicationbox}
The integral above is the backbone of probability and statistics. The area under the bell curve $e^{-x^2}$ (Figure~\ref{ch30:fig-gaussian}) between two points is what every ``what fraction of outcomes fall within so many standard deviations'' question reduces to, and no elementary formula for it exists. Statistical tables and software evaluate it by summing exactly the series of Example~\ref{ch30:ex-sub} (or a cleverer variant). The same situation arises in physics -- Fresnel integrals for light diffraction, the integral of $\sin(x)/x$ in optics -- where the only real way to get numbers is the power series. Series are not a curiosity; for a huge swath of applied mathematics they are the \emph{only} bridge from a non-elementary function to an actual decimal.
\end{applicationbox}

\begin{example}[Multiplication: the first terms of $e^x \sin x$]
\label{ch30:ex-mult}
To find the series of $e^x \sin x$ up to $x^3$, multiply the two series and collect by total degree:
\[
  e^x = 1 + x + \tfrac{x^2}{2} + \tfrac{x^3}{6} + \cdots,
  \qquad
  \sin x = x - \tfrac{x^3}{6} + \cdots.
\]
Collect powers:
\[
\begin{aligned}
  x^1:&\quad 1\cdot x = x,\\
  x^2:&\quad x\cdot x = x^2,\\
  x^3:&\quad \tfrac{x^2}{2}\cdot x + 1\cdot\left(-\tfrac{x^3}{6}\right) = \tfrac{x^3}{2} - \tfrac{x^3}{6} = \tfrac{x^3}{3}.
\end{aligned}
\]
So $e^x \sin x = x + x^2 + \tfrac{1}{3}x^3 + \cdots$. To get the same three terms from the Taylor definition you would differentiate $e^x \sin x$ three times with the product rule -- much more painful. Multiplying known series is the smart route.
\end{example}

\begin{example}[Building a Taylor series at a shifted center]
\label{ch30:ex-shift}
To expand $\ln x$ around $a = 1$ (you cannot center at $0$, since $\ln 0$ is undefined), write $x = 1 + (x-1)$ and reuse the known $\ln(1+u)$ series with $u = x-1$:
\[
  \ln x = \ln\bigl(1 + (x-1)\bigr)
  = (x-1) - \frac{(x-1)^2}{2} + \frac{(x-1)^3}{3} - \cdots,
  \qquad 0 < x \le 2.
\]
This is the Taylor series of $\ln x$ centered at $1$. The substitution did all the work the derivative computation would have done -- and the radius of convergence $R=1$ tells us the series is valid only out to $x=0$ and $x=2$, exactly up to the wall at $\ln 0$.
\end{example}

\begin{trapbox}
Series manipulation is forgiving, but three mistakes recur:
\begin{itemize}
  \item \textbf{Forgetting the new radius after substitution.} Substituting $u=-x^2$ into a series valid for $\abs{u}<1$ gives one valid for $\abs{x^2}<1$, i.e.\ $\abs{x}<1$ -- you must translate the condition.
  \item \textbf{Mismatched starting indices when adding.} Before combining $\sum$ and $\sum$, line up the powers; a series starting at $n=0$ and one starting at $n=1$ cannot be merged coefficient-by-coefficient until you peel off the stray term.
  \item \textbf{Assuming $\ln x$ has a Maclaurin series.} It does not -- $\ln 0$ is undefined, so you must shift the center, as in Example~\ref{ch30:ex-shift}.
\end{itemize}
\end{trapbox}

\section{The Catch: When Does the Series Equal the Function?}
\label{ch30:sec-catch}

We have been a little cavalier. The Taylor recipe produces a series from any infinitely differentiable function -- but does that series actually \emph{sum to} the function? For $e^x$, $\sin x$, $\cos x$, and all our examples, yes. But it is not automatic, and we should not paper over the gap.

\begin{historybox}
For a century after Newton and Leibniz, mathematicians used power series with breathtaking freedom and almost no rigor. Euler, the supreme calculator, manipulated divergent series, summed them, and routinely got correct and astonishing results -- but also occasionally absurd ones. The trouble was that ``$f(x)$ equals its Taylor series'' was treated as obvious. It took \textbf{Augustin-Louis Cauchy} in the 1820s to insist on a precise notion of convergence and to confront the uncomfortable question head-on. Cauchy himself produced the famous warning example below. The full clean-up -- a rigorous theory of when a Taylor series converges \emph{to its function}, via the remainder term -- is the subject of the next chapter, Taylor's theorem.
\end{historybox}

\begin{trapbox}[Cauchy's cautionary function]
Consider
\[
  f(x) = \begin{cases} e^{-1/x^2}, & x \neq 0,\\[2pt] 0, & x = 0. \end{cases}
\]
This function is infinitely differentiable everywhere, and \emph{every one} of its derivatives at $0$ equals zero: $f(0) = f'(0) = f''(0) = \cdots = 0$. So its Maclaurin series is
\[
  0 + 0\cdot x + 0\cdot x^2 + \cdots = 0,
\]
the zero series. But $f(x)$ is plainly \emph{not} zero for $x\neq 0$. Here is a perfectly smooth function whose Taylor series converges everywhere -- to the wrong answer. The series sees only what happens at the single point $x=0$, and at that point $f$ is flatter than every polynomial. The function hides all its behavior ``infinitely far'' from the origin in the language of derivatives.
\end{trapbox}

\begin{intuitionbox}
The moral: a Taylor series is built from purely \emph{local} data -- the derivatives at one point. Sometimes (most happy times) that local fingerprint determines the function globally, and the series reproduces it perfectly; such functions are called \emph{analytic}. Sometimes, as with Cauchy's example, the fingerprint is misleading. The tool that decides which case you are in is the \emph{remainder}: the difference $f(x) - P_N(x)$ between the function and its Taylor polynomial. The series equals the function exactly when this remainder shrinks to $0$ as $N\to\infty$. Proving the remainder vanishes -- for $e^x$, $\sin x$, $\cos x$, and the rest -- is the precise, satisfying business of the next chapter.
\end{intuitionbox}

\begin{aestheticbox}
Step back and admire what we have found. The transcendental functions -- exponential growth, oscillation, the logarithm, the inverse trigonometric functions -- are not a wild zoo of unrelated objects. Each is, within its interval, a single infinite polynomial with a clean, patterned list of coefficients. Differentiation, integration, multiplication, and substitution move among them as fluidly as among ordinary polynomials. The factorials, the alternating signs, the odd-and-even symmetries -- all of it is bookkeeping for one idea: \emph{match a function to a polynomial at a point, to every order at once.} That a smooth function should be ``secretly a polynomial of infinite degree'' is, after limits themselves, perhaps the most beautiful structural truth in calculus -- and the one that makes the subject \emph{computable}.
\end{aestheticbox}

\section{A Table to Carry With You}
\label{ch30:sec-table}

These are the series worth memorizing. Everything else is built from them.
\[
\begin{array}{llll}
\dfrac{1}{1-x} &=\ \displaystyle\sum_{n=0}^{\infty} x^n &=\ 1 + x + x^2 + x^3 + \cdots & \abs{x}<1 \\[1.6em]
e^x &=\ \displaystyle\sum_{n=0}^{\infty} \dfrac{x^n}{n!} &=\ 1 + x + \dfrac{x^2}{2!} + \dfrac{x^3}{3!} + \cdots & \text{all } x \\[1.6em]
\sin x &=\ \displaystyle\sum_{n=0}^{\infty} \dfrac{(-1)^n x^{2n+1}}{(2n+1)!} &=\ x - \dfrac{x^3}{3!} + \dfrac{x^5}{5!} - \cdots & \text{all } x \\[1.6em]
\cos x &=\ \displaystyle\sum_{n=0}^{\infty} \dfrac{(-1)^n x^{2n}}{(2n)!} &=\ 1 - \dfrac{x^2}{2!} + \dfrac{x^4}{4!} - \cdots & \text{all } x \\[1.6em]
\ln(1+x) &=\ \displaystyle\sum_{n=1}^{\infty} \dfrac{(-1)^{n-1} x^n}{n} &=\ x - \dfrac{x^2}{2} + \dfrac{x^3}{3} - \cdots & -1 < x \le 1 \\[1.6em]
\arctan x &=\ \displaystyle\sum_{n=0}^{\infty} \dfrac{(-1)^n x^{2n+1}}{2n+1} &=\ x - \dfrac{x^3}{3} + \dfrac{x^5}{5} - \cdots & -1 \le x \le 1
\end{array}
\]

\section{Exercises}
\label{ch30:sec-exercises}

\begin{exercise}[Radius and interval of convergence]
For each series, find the radius of convergence, then test the endpoints to determine the full interval of convergence.
\[
  \text{(a) } \sum_{n=1}^{\infty}\frac{x^n}{n^2},
  \qquad
  \text{(b) } \sum_{n=0}^{\infty}\frac{x^n}{2^n},
  \qquad
  \text{(c) } \sum_{n=1}^{\infty}\frac{(x-3)^n}{n\,5^n},
  \qquad
  \text{(d) } \sum_{n=0}^{\infty} (2n)!\,x^n.
\]
\end{exercise}

\begin{exercise}[Series by manipulation]
Without computing a single derivative, find the Maclaurin series (and its radius of convergence) for:
\[
  \text{(a) } \frac{1}{1+3x},
  \qquad
  \text{(b) } x^2 e^{-x},
  \qquad
  \text{(c) } \cos(2x),
  \qquad
  \text{(d) } \frac{x}{1-x^2}.
\]
\end{exercise}

\begin{exercise}[Differentiate a known series]
Starting from $\dfrac{1}{1-x} = \sum_{n=0}^\infty x^n$, differentiate twice to find a closed-form sum for $\sum_{n=2}^{\infty} n(n-1)x^{n-2}$, valid for $\abs{x}<1$. Then evaluate $\sum_{n=2}^{\infty} n(n-1)(1/2)^{n-2}$ as an exact number.
\end{exercise}

\begin{exercise}[Integrate a known series]
Use the series for $\dfrac{1}{1+x^2}$ to write $\displaystyle\int_0^{1/2}\frac{\dd x}{1+x^2}$ as an infinite series, and use the first three terms to estimate $\arctan(1/2)$. Compare with the true value.
\end{exercise}

\begin{exercise}[Computing a transcendental number]
Use the Maclaurin series for $e^x$ at $x=1$ to estimate $e$ using the terms up to $1/5!$. How many decimal places do you trust, and why does the alternating-series-style reasoning \emph{not} directly apply here (the terms do not alternate)? Bound the tail $\sum_{n\ge 6} 1/n!$ by comparison with a geometric series.
\end{exercise}

\begin{exercise}[A limit by series]
Evaluate $\displaystyle\lim_{x\to 0}\frac{\sin x - x}{x^3}$ by replacing $\sin x$ with the first few terms of its series, simplifying, and taking the limit. Then confirm with L'Hopital's rule, and reflect on which method exposed \emph{why} the answer is what it is.
\end{exercise}

\begin{exercise}[Product of series]
Find the Maclaurin series of $\dfrac{e^x}{1-x}$ up to and including the $x^3$ term, by multiplying the series for $e^x$ and $\dfrac{1}{1-x}$. (You should find the coefficient of $x^n$ is $\sum_{k=0}^{n} 1/k!$.)
\end{exercise}

\begin{exercise}[Shifted center]
Find the Taylor series of $f(x) = 1/x$ centered at $a = 2$. (Hint: write $\dfrac{1}{x} = \dfrac{1}{2 + (x-2)} = \dfrac{1}{2}\cdot\dfrac{1}{1 + (x-2)/2}$ and use the geometric series.) State the radius of convergence and explain geometrically why it is what it is.
\end{exercise}

\begin{exercise}[Conceptual: uniqueness]
Explain in your own words why a function can have only one power series at a given center. If $\sum a_n x^n = \sum b_n x^n$ for all $x$ in some interval around $0$, prove that $a_n = b_n$ for every $n$. (Hint: differentiate $n$ times and set $x=0$.)
\end{exercise}

\begin{exercise}[Proof: term-by-term and the derivative]
Let $f(x) = \sum_{n=0}^{\infty} c_n x^n$ have radius $R>0$, and let $g(x) = \sum_{n=1}^{\infty} n c_n x^{n-1}$ be the formally differentiated series. Prove that $g$ has the same radius of convergence $R$ as $f$. (You do not need to prove $g = f'$; just the radius. Use the ratio-test characterization and the fact that $\sqrt[n]{n}\to 1$ or that $n/(n+1)\to 1$.)
\end{exercise}

\begin{exercise}[Cauchy's function, conceptual]
For $f(x) = e^{-1/x^2}$ (with $f(0)=0$), explain why every derivative at $0$ is zero, at least intuitively, by arguing that $e^{-1/x^2}$ goes to $0$ as $x\to 0$ faster than any power of $x$. What does this example tell you about the difference between ``the Taylor series converges'' and ``the Taylor series converges \emph{to the function}''?
\end{exercise}

\begin{exercise}[Euler's formula by series]
Substitute $x = i\theta$ into the series for $e^x$, group the real and imaginary terms separately, and recognize the series for $\cos\theta$ and $\sin\theta$. Conclude $e^{i\theta} = \cos\theta + i\sin\theta$, and deduce the value of $e^{i\pi}$.
\end{exercise}
